\section{G. 数值方法与计算数学}

\subsection{G1. 常微分方程数值解}

\subsubsection{经典 RK4}

\begin{align*}
k_1 &= f(t_n, y_n) \\
k_2 &= f(t_n + h/2,\; y_n + hk_1/2) \\
k_3 &= f(t_n + h/2,\; y_n + hk_2/2) \\
k_4 &= f(t_n + h,\; y_n + hk_3) \\
y_{n+1} &= y_n + \frac{h}{6}(k_1 + 2k_2 + 2k_3 + k_4)
\end{align*}

\subsubsection{辛积分器}

\textbf{Verlet 积分：}
\[
x_{n+1} = 2x_n - x_{n-1} + a_n h^2
\]

\textbf{Leapfrog（蛙跳）：}
\begin{align*}
v_{n+1/2} &= v_{n-1/2} + a_n h \\
x_{n+1} &= x_n + v_{n+1/2} h
\end{align*}

辛积分器能量在 bounded 范围内振荡，不会系统性漂移。布料模拟、刚体模拟、粒子系统都用辛积分器。

\subsection{G2. 偏微分方程数值解}

\subsubsection{有限差分法（FDM）}

\[
\frac{\partial u}{\partial x} \approx \frac{u_{i+1} - u_{i-1}}{2h}, \quad
\frac{\partial^2 u}{\partial x^2} \approx \frac{u_{i+1} - 2u_i + u_{i-1}}{h^2}
\]

\textbf{CFL 条件：} 对波动方程，$\dfrac{c\,\Delta t}{\Delta x} \leq 1$。

\subsubsection{有限元法（FEM）}

\textbf{弱形式：}
\[
\int_\Omega \nabla u \cdot \nabla v\,d\Omega = \int_\Omega fv\,d\Omega
\]

\textbf{刚度矩阵：} $K_{ij} = \displaystyle\int_\Omega \nabla\phi_i \cdot \nabla\phi_j\,d\Omega$

\subsubsection{谱方法}

用全局正交基（傅里叶、切比雪夫、球谐）展开解。对光滑解指数收敛。

\subsection{G3. 数值线性代数}

\begin{center}
\begin{tabular}{lll}
\toprule
\textbf{分解} & \textbf{形式} & \textbf{应用} \\
\midrule
LU & $A = LU$ & 解线性系统 \\
QR & $A = QR$ & 最小二乘、特征值 \\
SVD & $A = U\Sigma V^T$ & PCA、伪逆、低秩近似 \\
Cholesky & $A = LL^T$（正定） & 快速求解 \\
特征分解 & $A = P\Lambda P^{-1}$ & 谱分析、矩阵函数 \\
Schur & $A = QTQ^*$ & 数值稳定的特征值计算 \\
\bottomrule
\end{tabular}
\end{center}

\textbf{迭代法：} 共轭梯度法（CG）、GMRES、多重网格。

\subsection{G4. 逼近论}

\subsubsection{Padé 近似}

\[
f(x) \approx \frac{P_m(x)}{Q_n(x)} = \frac{a_0 + a_1 x + \cdots + a_m x^m}{1 + b_1 x + \cdots + b_n x^n}
\]

可以逼近有极点的函数（泰勒级数做不到），收敛半径通常更大。

\subsubsection{切比雪夫逼近}

切比雪夫节点：$x_k = \cos\dfrac{(2k+1)\pi}{2(n+1)}$（避免 Runge 现象）。

\subsection{G5. 优化}

\begin{itemize}
  \item 梯度下降：$x_{k+1} = x_k - \alpha \nabla f(x_k)$
  \item 牛顿法：$x_{k+1} = x_k - H^{-1}\nabla f$
  \item L-BFGS（拟牛顿，CG 中广泛使用）
  \item 拉格朗日乘子法、KKT 条件
\end{itemize}